% !TEX TS-program = lualatex

% ------------------------------------------- %
% -- AUTOMATICALLY GENERATED - DO NOT EDIT -- %
% ------------------------------------------- %

\documentclass{tutodoc}

\usepackage{../preamble.cfg}


\begin{document}

\subsection{CSS}
\label{products-css}

\begin{tdocnote}
The {\html}, {\css}, and {\javascript} files for product development and demonstration were created using {\claudeAI} and {\geminiAI} AI assistants.
\end{tdocnote}

Each palette color is defined as an individual variable named according to the pattern \tdoccodein{css}+--pal<name>-<nb>+, where \tdoccodein{css}+<name>+ is the standard palette name and \tdoccodein{css}+<nb>+ is the desired index ranging. Each palette color variable is defined as an \tdoccodein{text}+RGB+ value using percentage notation. For example, the file \tdoccodein{text}+palettes-hf/Accent.css+ looks like the following partial code.

\begin{tdoccode}{css}
:root {
  --palAccent-1: rgb(49.803922% 78.823529% 49.803922%);
  --palAccent-2: rgb(74.509804% 68.235294% 83.137255%);
  --palAccent-3: rgb(99.215686% 75.294118% 52.54902%);
  /* Other RBG colors.*/
}
\end{tdoccode}

The following example illustrates how to generate gradient variables via selective color extraction and custom reordering, while using a standalone color for warning text.

\begin{tdoccode}{css}
:root {
  --transformed-accent-gradient: linear-gradient(
    90deg,
    var(--palAccent-6),
    var(--palAccent-3),
    var(--palAccent-8),
    var(--palAccent-1)
  );
}

.transformed-accent-gradient {
  background: var(--transformed-accent-gradient);
}

.warning-text {
  color: var(--palAccent-3);
}
\end{tdoccode}

\end{document}
